\documentclass[11pt]{article}
\usepackage[margin=1in]{geometry}
\usepackage{amsmath,amssymb,amsthm,mathtools}
\usepackage[hidelinks]{hyperref}
\usepackage{microtype}

\newtheorem{theorem}{Theorem}[section]
\newtheorem{proposition}[theorem]{Proposition}
\newtheorem{lemma}[theorem]{Lemma}
\newtheorem{corollary}[theorem]{Corollary}
\theoremstyle{definition}
\newtheorem{definition}[theorem]{Definition}
\theoremstyle{remark}
\newtheorem{remark}[theorem]{Remark}

\newcommand{\Z}{\mathbb Z}
\newcommand{\F}{\mathbb F}
\newcommand{\C}{\mathbb C}
\newcommand{\U}{\mathcal U}
\newcommand{\Ad}{\operatorname{Ad}}
\newcommand{\tr}{\operatorname{tr}}
\newcommand{\Cl}{\operatorname{Cl}}
\newcommand{\Fix}{\operatorname{Fix}}
\newcommand{\dist}{\operatorname{dist}}
\newcommand{\rank}{\operatorname{rank}}
\newcommand{\Tr}{\operatorname{Tr}}

\title{An Eight-Lamp Spin Reduction and Its Regular-Tensor Boundary}
\author{Research note derived from the marked group in
  \texttt{non\_mf\_groups\_exist.tex}}
\date{13 August 2026}

\begin{document}
\maketitle

\begin{abstract}
We isolate an exact finite-dimensional structure forced by the marked
central involution in the group $E$ of
\texttt{non\_mf\_groups\_exist.tex}.  Write
$\Gamma=\Z^3\rtimes \mathrm{SL}_3(\Z)$, let
$\alpha(v,A)=(2v,A)$, put $H=\alpha(\Gamma)$, and set
$d=tct^{-1}$ and $w=[d,ada^{-1}]$.  In every tracial representation in
which $w=-1$, the eight conjugates of $d$ indexed by $\Gamma/H$ pairwise
anticommute.  They generate a canonical copy of
$\Cl_8(\C)\cong M_{16}(\C)$, and the entire $\Gamma$-action splits as a
fixed projective spin action tensored with an inverse-projective coefficient
action.  In a hypothetical hyperlinear embedding the negative central
corner has trace $1/2$.

We then prove a limitation which is not visible from the abstract spin
normal form.  In every hypothetical trace-preserving embedding, the
normalized negative-corner trace restricts to the regular character on both
$\Gamma$ and $H$.  Consequently the normalized characters of the coordinate
adjoint representations converge to the regular character, and the
coordinate low-spectral spaces have vanishing normalized dimension.  This
remains true even though they contain the distinguished unit-norm lamp
vector.  Tensoring with finite regular representations of an explicit
linear quotient preserves faithfulness, the negative corner, and every
Clifford statistic, showing that no hidden choice of microstates repairs the
rank loss.  The fixed $M_{16}$ factor therefore does not yield a
positive-density adjoint obstruction.  Any completion of this route must
use a relative or center-valued multiplicity invariant compatible with the
HNN letter, or a simultaneous coordinate-commutant recovery theorem stronger
than ordinary spectral-gap recovery.

This paper is a proved reduction and a proved stopping theorem, not a proof
of an explicit nonhyperlinear group.  The existence of a nonhyperlinear
discrete group remains open in the primary literature cited below.
\end{abstract}

\section{Status}

The existence of a nonhyperlinear discrete group remains open.  Dogon and
Vigdorovich obtain conditional nonhyperlinear finite central extensions from
Hilbert--Schmidt stability and isolate a concrete commensurator-stability
question \cite{DogonVigdorovich}.  Taller and Vidick prove RE-hardness for
linear-constraint-system games with completeness $1-\varepsilon$ and state
that perfect completeness would imply a nonhyperlinear group
\cite{TallerVidick}.  Alekseev and Thom prove coordinate centralizer
recovery for sofic approximations of Kazhdan groups, while formulating the
tracial-matrix analogue as Open Problem~6.2 \cite{AlekseevThom}.

Accordingly, no missing matrix-coordinate implication will be used below.
The purpose of the note is twofold: to record the exact spin factor forced by
the marked presentation, and to prove that ordinary normalized adjoint rank
cannot close the resulting coefficient problem.

\section{Tracial coordinate conventions}

Fix a free ultrafilter $\omega$ on $\mathbb N$.  For matrix sizes $m_n$ put
\[
 \mathcal M_\omega
 =\prod_\omega(M_{m_n}(\C),\tr_{m_n}),
 \qquad
 \|x\|_{2,n}=\tr_{m_n}(x^*x)^{1/2}.
\]
Here the tracial ultraproduct is the quotient of the bounded product by the
ideal of sequences whose normalized $2$-norm tends to zero along $\omega$.
A trace-preserving embedding of $L(E)$ means a unital normal
$*$-monomorphism $\theta:L(E)\to\mathcal M_\omega$ preserving the canonical
traces.  Equivalently, its group-unitary lifts form hyperlinear microstates
whose traces converge to the regular character.

We will repeatedly use the following elementary lifting fact.

\begin{lemma}[Corners and unitary lifts]\label{lem:cornerlifts}
Let $p\in\mathcal M_\omega$ be a projection of trace $s>0$.  There are
projection representatives $p_n\in M_{m_n}(\C)$, with
$r_n:=\rank(p_n)$ and $r_n/m_n\to_\omega s$, for which
\[
 p\mathcal M_\omega p
 \cong \prod_\omega(M_{r_n}(\C),\tr_{r_n})
\]
trace preservingly.  Every unitary in the corner has representatives which
are unitaries in the matrix corners $p_nM_{m_n}(\C)p_n$.
\end{lemma}

\begin{proof}
Lift $p$ by bounded matrices and replace the lifts by the spectral
projections for the interval $[1/2,\infty)$.  The resulting sequence still
represents $p$, and normalized trace convergence gives
$r_n/m_n\to_\omega s$.  Compression identifies the corner quotient with the
displayed tracial ultraproduct after renormalizing by $r_n$.  Finally, a
lift of a corner unitary is asymptotically unitary in normalized $2$-norm;
its polar part, completed arbitrarily on its kernel, is a unitary lift at
$2$-norm distance tending to zero.
\end{proof}

\begin{lemma}[Exact coordinate lifting of a matrix packet]
\label{lem:matrixlift}
Let $B\cong M_k(\C)$ be a unital subalgebra of
$p\mathcal M_\omega p$.  After changing the projection representatives in
Lemma~\ref{lem:cornerlifts} by normalized $2$-norm zero, one may arrange
that $r_n=k\ell_n$ and choose unital embeddings
\[
 \iota_n:M_k(\C)\longrightarrow M_{r_n}(\C)
\]
whose ultraproduct is the given inclusion of $B$.  After unitary coordinate
changes, $\iota_n(M_k)=M_k(\C)\otimes1_{\ell_n}$.
\end{lemma}

\begin{proof}
Choose matrix units $(e_{ij})$ for $B$ and bounded coordinate lifts.  Their
matrix-unit relations hold in normalized $2$-norm along $\omega$.
Spectral truncation of the diagonal lifts, followed by successive
orthogonalization, gives pairwise orthogonal projections representing the
$e_{ii}$.  The polar parts of the compressed lifts of $e_{i1}$ give partial
isometries between subprojections whose missing normalized traces tend to
zero.  Cutting all initial and final projections to their common minimum
rank produces exact partial isometries $v_{i,n}$ with a common initial
projection.  Then
\[
 f_{ij,n}=v_{i,n}v_{j,n}^*
\]
are exact matrix units and still represent $e_{ij}$.  Their support
$p'_n=\sum_i f_{ii,n}$ represents $p$ and has rank $k\ell_n$.  Replacing
$p_n$ by $p'_n$ proves the first assertion.  The classification of unital
representations of a full matrix algebra gives the displayed tensor form.
\end{proof}

\section{The marked group and its affine quotient}

Let
\[
 \Gamma=\Z^3\rtimes \mathrm{SL}_3(\Z),
 \qquad (v,A)(u,B)=(v+Au,AB),
\]
and define
\[
 \alpha(v,A)=(2v,A),\qquad H=\alpha(\Gamma),
 \qquad a=(e_1,1).
\]
Let $E$ be the finitely presented group obtained by adjoining $t,c$ with
\[
 t\gamma t^{-1}=\alpha(\gamma)\quad(\gamma\in\Gamma),
 \qquad c^2=1,\qquad[c,\Gamma]=1,
\]
and imposing that
\[
 d:=tct^{-1},\qquad w:=[d,ada^{-1}]
\]
is a central involution.  The companion non-MF paper proves $w\ne1$ in
$E$ by an explicit infinite Clifford representation.

The groups $\Gamma$ and $H\cong\Gamma$ are infinite and have property
\textup{(T)} \cite[Example~1.7.4(i)]{BHV}.  We use the commutator convention
$[x,y]=xyx^{-1}y^{-1}$.

\begin{lemma}[Affine two-transitivity]\label{lem:2trans}
The left action of $\Gamma$ on $\Omega:=\Gamma/H$ is permutation-isomorphic
to the affine action of
\[
 \F_2^3\rtimes \mathrm{GL}_3(\F_2)=\mathrm{AGL}_3(2)
\]
on $\F_2^3$.  In particular, $|\Omega|=8$ and the action is
$2$-transitive.
\end{lemma}

\begin{proof}
Reduce the translation coordinate modulo $2$.  The element $(v,A)$ acts on
$x\in\F_2^3$ by $x\mapsto \bar v+\bar A x$.  Since
$\mathrm{SL}_3(\F_2)=\mathrm{GL}_3(\F_2)$, the stabilizer of $0$ consists
of the elements whose translation coordinate is even, namely $H$.
The stabilizer $\mathrm{GL}_3(\F_2)$ is transitive on the seven nonzero
vectors, proving two-transitivity.
\end{proof}

\begin{lemma}[Linear quotient and separation]\label{lem:linearquotient}
Embed $\Gamma$ in $\mathrm{GL}_4(\Z)$ by
\[
 (v,A)\longmapsto
 \begin{pmatrix}A&v\\0&1\end{pmatrix},
 \qquad
 D=\operatorname{diag}(2,2,2,1),
\]
and put $\bar G=\langle\Gamma,D\rangle\le\mathrm{GL}_4(\Z[1/2])$.  The assignments
\[
 q|_\Gamma=\operatorname{id}_\Gamma,\qquad q(t)=D,\qquad q(c)=1
\]
define a homomorphism $q:E\to\bar G$.  It satisfies $q(w)=1$, is injective
on $\Gamma$, and
\[
 \langle w\rangle\cap\Gamma=\{1\}.
\]
Moreover, $\bar G$ is residually finite.
\end{lemma}

\begin{proof}
Direct matrix multiplication gives
\[
 D\begin{pmatrix}A&v\\0&1\end{pmatrix}D^{-1}
 =\begin{pmatrix}A&2v\\0&1\end{pmatrix},
\]
so the HNN relations are respected.  Every relation involving $c$ or $w$
maps to the identity, proving that $q$ is well defined and that $q(w)=1$.
The displayed affine embedding of $\Gamma$ is faithful.  If the nontrivial
involution $w$ belonged to $\Gamma$, injectivity of $q|_\Gamma$ would
contradict $q(w)=1$; this proves the intersection assertion.  Finally,
$\bar G$ is a finitely generated linear group in characteristic zero, so it
is residually finite by Mal'cev's theorem \cite{Malcev}.
\end{proof}

For $x=rH\in\Omega$, define $d_x=rdr^{-1}$.  This is well defined because
$d$ centralizes $H$: if $h=t\gamma t^{-1}\in H$, then
\[
 d h d^{-1}=t(c\gamma c^{-1})t^{-1}=h.
\]

\section{The exact Clifford packet}

\begin{theorem}[Spin-factor extraction]\label{thm:spin}
Let $\pi:E\to\U(M)$ be a homomorphism into a finite tracial von Neumann
algebra and put
\[
 p=\frac{1-\pi(w)}2.
\]
If $p\ne0$, compress to $pMp$ and normalize its trace.  Then:
\begin{enumerate}
\item $\pi(w)=-1$ on the compressed algebra;
\item the eight self-adjoint involutions
      $D_x:=\pi(d_x)$, $x\in\Omega$, satisfy
      \[
       D_xD_y=-D_yD_x\qquad(x\ne y);
      \]
\item $A:=W^*(D_x:x\in\Omega)$ is a full matrix factor
      isomorphic to $M_{16}(\C)$;
\item $\pi(\Gamma)$ normalizes $A$.
\end{enumerate}
\end{theorem}

\begin{proof}
Centrality of $w$ makes $p$ commute with $\pi(E)$, and the compression has
$\pi(w)=-1$.  The pair $(H,aH)$ consists of two distinct points of
$\Omega$.  By Lemma~\ref{lem:2trans}, for every ordered pair $x\ne y$
there is $g\in\Gamma$ sending $(H,aH)$ to $(x,y)$.  Centrality of $w$
therefore gives
\[
 [D_x,D_y]
 =\pi\bigl(g[d,ada^{-1}]g^{-1}\bigr)
 =\pi(w)=-1.
\]
For involutions this is equivalent to anticommutation.

The universal complex $C^*$-algebra on eight self-adjoint involutions with
these relations is $\Cl_8(\C)\cong M_{16}(\C)$.  The resulting unital
$*$-homomorphism into $pMp$ is injective because $M_{16}(\C)$ is simple.
Finally, conjugation by $g\in\Gamma$ sends $D_x$ to $D_{gx}$, so it
normalizes their matrix algebra.
\end{proof}

\begin{corollary}[Coefficient normal form]\label{cor:normal}
Under the hypotheses of Theorem~\ref{thm:spin}, there is a finite tracial
von Neumann algebra $N$ and an isomorphism
\[
 pMp\cong M_{16}(\C)\,\bar\otimes\,N
\]
which carries $A$ onto $M_{16}(\C)\otimes1$.  There are unitaries
$S_g\in U(16)$ and $v_g\in\U(N)$ such that
\[
 \pi(g)=S_g\otimes v_g\qquad(g\in\Gamma).
\]
The $S_g$ implement the affine permutation action on the eight Clifford
generators.  Thus
\[
 S_gS_h=\mu(g,h)S_{gh},\qquad
 v_gv_h=\mu(g,h)^{-1}v_{gh}
\]
for a scalar multiplier $\mu$, and the projective action
$g\mapsto\Ad S_g$ factors through $\mathrm{AGL}_3(2)$.
\end{corollary}

\begin{proof}
The standard matrix-unit decomposition for a unital full matrix subfactor
gives $pMp\cong M_{16}(\C)\bar\otimes N$.  Since $\pi(g)$ normalizes the
first tensor factor and every automorphism of $M_{16}(\C)$ is inner, choose
$S_g$ implementing that automorphism.  Then
$(S_g^*\otimes1)\pi(g)$ commutes with the first factor and hence equals
$1\otimes v_g$.  Multiplicativity of $\pi$ gives the two inverse cocycle
identities.  The action on the Clifford generators is the action on
$\Gamma/H$, so its projectivization factors through the affine quotient.
\end{proof}

\begin{corollary}[Exact positive-density coordinate spin blocks]
\label{cor:coordinatespin}
Suppose that $\theta:L(E)\to\mathcal M_\omega$ is a trace-preserving
embedding and put $p=(1-\theta(w))/2$.  One can represent $p$ by projections
$p_n$ of ranks $r_n=16k_n$ so that
\[
 \frac{r_n}{m_n}\longrightarrow_\omega\frac12
\]
and represent the Clifford algebra of Theorem~\ref{thm:spin} by exact
coordinate inclusions
\[
 M_{16}(\C)\otimes1_{k_n}\subseteq M_{16k_n}(\C)=p_nM_{m_n}(\C)p_n.
\]
Equivalently, the eight group-unitary representatives may be changed by
normalized $2$-norm zero to exact self-adjoint involutions which pairwise
anticommute at every coordinate.
In particular, the coefficient commutant has vector-space dimension $k_n^2$,
exactly $1/256$ of the vector-space dimension of the corner matrix algebra.
\end{corollary}

\begin{proof}
The regular trace gives $\tau(\theta(w))=0$, so $\tau(p)=1/2$.  Apply
Lemma~\ref{lem:matrixlift} with $k=16$ to the copy of $M_{16}(\C)$ from
Theorem~\ref{thm:spin}.  The commutant of
$M_{16}(\C)\otimes1_{k_n}$ in $M_{16k_n}(\C)$ is
$1_{16}\otimes M_{k_n}(\C)$, which has the asserted dimension ratio.
\end{proof}

\begin{remark}[What is, and is not, positive density]\label{rem:density}
Corollary~\ref{cor:coordinatespin} makes the positive-density assertion
literal at matrix coordinates: the negative central corner has asymptotic
trace $1/2$, and its Clifford coefficient commutant occupies exactly a
$1/256$ fraction of the corner's matrix-vector space.

This does not imply that the $H$-fixed part of that coefficient commutant,
or the excess $H$-fixed part over the $\Gamma$-fixed part, has positive
normalized dimension.  The coefficient action $v_g$ can carry essentially
all of the adjoint spectral mass.  Section~\ref{sec:dilution} proves that
ordinary adjoint density is already absent for trace-preserving lifts and
remains absent under harmless regular tensoring.
\end{remark}

\begin{proposition}[Regular adjoint character in the negative corner]
\label{prop:cornerregular}
Suppose that $\theta:L(E)\to\mathcal M_\omega$ is a trace-preserving
embedding, put $p=(1-\theta(w))/2$, and equip $p\mathcal M_\omega p$ with
its normalized trace.  The restrictions of this trace to both $\Gamma$ and
$H$ are their regular characters.

More precisely, use Lemma~\ref{lem:cornerlifts} and let $U_{g,n}\in U(r_n)$
be unitary representatives of $p\theta(g)$.  For
$K\in\{\Gamma,H\}$ choose a finite symmetric generating set $S_K$ and put
\[
 \Delta_{K,n}
 =\frac1{2|S_K|}\sum_{s\in S_K}
   (1-\Ad U_{s,n})^*(1-\Ad U_{s,n})
 \quad\text{on }L^2(M_{r_n},\tr_{r_n}).
\]
If $\Delta_K^{\mathrm{reg}}$ is the same Laplacian in the left regular
representation and
\[
 \kappa_K=\inf\bigl(\sigma(\Delta_K^{\mathrm{reg}})\bigr)>0,
\]
then, for every $0<\eta<\kappa_K$,
\begin{align*}
 \tr_{r_n^2}(\Ad U_{g,n})&\longrightarrow_\omega\delta_e(g)
   &&(g\in K),\\
 \frac{\rank\,1_{[0,\eta]}(\Delta_{K,n})}{r_n^2}
   &\longrightarrow_\omega0.&
\end{align*}
\end{proposition}

\begin{proof}
Since $w\ne1$ and the embedding is trace preserving, $\tau(\theta(w))=0$
and $\tau(p)=1/2$.  For $g\in\Gamma$ the normalized corner trace is
\[
 \tau_p(p\theta(g))
 =\tau(\theta(g))-\tau(\theta(wg)).
\]
If $g=e$, this equals $1$.  If $g\ne e$, then both $g$ and $wg$ are
nonidentity elements of $E$: the second assertion uses
Lemma~\ref{lem:linearquotient}, which gives $w\notin\Gamma$.  Both trace
terms therefore vanish.  Thus the restriction is $\delta_e$, and the same
calculation applies to $H\le\Gamma$.

If $U_{g,n}$ are compressed unitary representatives in matrices of size
$r_n$, the normalized character of their conjugation action is
\[
 \tr_{r_n^2}(\Ad U_{g,n})
 =\left|\tr_{r_n}(U_{g,n})\right|^2.
\]
It therefore converges pointwise to $\delta_e$ for either group.  Expanding
a power of $\Delta_{K,n}$ gives finitely many products of coordinate
conjugations.  The unitary lifts are multiplicative modulo normalized
$2$-norm zero, so the preceding character convergence gives convergence of
every moment of $\Delta_{K,n}$ to the corresponding moment of
$\Delta_K^{\mathrm{reg}}$.

The groups $\Gamma$ and $H\cong\Gamma$ are infinite property-(T) groups.
Thus their regular representations have no invariant vector and the
displayed numbers $\kappa_K$ are positive.  Given $0<\eta<\kappa_K$, choose
a continuous function $f:[0,2]\to[0,1]$ which is one on $[0,\eta]$ and
zero on $[\kappa_K,2]$.  Polynomial approximation and moment convergence
give
\[
 \frac{1}{r_n^2}\Tr f(\Delta_{K,n})\longrightarrow_\omega0.
\]
Since $1_{[0,\eta]}\le f$, the asserted rank convergence follows.
\end{proof}

\begin{corollary}[A unit lamp in an asymptotically zero-rank sector]
\label{cor:lampzerorank}
Under the hypotheses of Proposition~\ref{prop:cornerregular}, fix
$0<\eta<\kappa_H$ and let
$P_{H,n}=1_{[0,\eta]}(\Delta_{H,n})$.  There are unitary representatives
$D_n$ of $p\theta(d)$ and $A_n$ of $p\theta(a)$ such that
\begin{align*}
 \frac{\rank(P_{H,n})}{r_n^2}&\longrightarrow_\omega0,\\
 \|P_{H,n}D_n-D_n\|_{2,n}&\longrightarrow_\omega0,\\
 \|D_n-A_nD_nA_n^*\|_{2,n}^2&\longrightarrow_\omega2.
\end{align*}
Thus the exact $M_{16}$ packet supplies a distinguished asymptotically
$H$-fixed unit vector with an order-one Clifford wall, but supplies no
positive lower bound for the ordinary normalized dimension of the
$H$-low-spectral space.
\end{corollary}

\begin{proof}
The rank assertion is Proposition~\ref{prop:cornerregular}.  Since $d$
centralizes $H$, its ultraproduct vector is fixed by every
$\Ad(p\theta(h))$.  Hence
\[
 \langle\Delta_{H,n}D_n,D_n\rangle\longrightarrow_\omega0.
\]
The spectral theorem gives
\[
 \|(1-P_{H,n})D_n\|_{2,n}^2
 \le\eta^{-1}\langle\Delta_{H,n}D_n,D_n\rangle,
\]
which proves the second assertion.  In the negative corner, $d$ and
$ada^{-1}$ anticommute by Theorem~\ref{thm:spin}.  Traciality then gives
$\tau_p(dada^{-1})=0$, and consequently
\[
 \|d-ada^{-1}\|_2^2=2.
\]
Passing to the chosen representatives proves the last limit.
\end{proof}

\section{The coefficient commutant inclusion}

Continue in the negative corner and put
\[
 C=\pi(\Gamma)'\cap pMp,
 \qquad D=\pi(H)'\cap pMp.
\]
Since $H\le\Gamma$ and $H=t\Gamma t^{-1}$ in $E$,
\begin{equation}\label{eq:CD}
 C\subseteq D,
 \qquad D=\pi(t)C\pi(t)^*.
\end{equation}

\begin{proposition}[Eight-coset expectation]\label{prop:expect}
Choose left-coset representatives $r_1=e,r_2,\ldots,r_8$ for $\Gamma/H$.
Then
\[
 \mathsf E_C(x)=\frac18\sum_{i=1}^8
 \pi(r_i)x\pi(r_i)^*,\qquad x\in D,
\]
is the trace-preserving conditional expectation $D\to C$.  For $x\in D_+$,
\[
 \mathsf E_C(x)\ge\frac18x.
\]
Thus $C\subseteq D$ has Pimsner--Popa index at most $8$.
If $\pi(w)=-1$, then $d\in D\setminus C$, so the inclusion is proper.
\end{proposition}

\begin{proof}
For $x\in D$, changing a representative on the right by an element of $H$
does not change the corresponding conjugate.  Left multiplication by an
element of $\Gamma$ permutes the left cosets, so the average lies in $C$.
It fixes $C$ pointwise and preserves the trace, and hence is the conditional
expectation.  Positivity and the representative $r_1=e$ give the displayed
lower bound.  The element $d$ centralizes $H$, hence lies in $D$.  If it
also lay in $C$, then it would commute with $a\in\Gamma$, forcing
$[d,ada^{-1}]=1$, contrary to $\pi(w)=-1$.
\end{proof}

In finite-dimensional representation spaces, the inclusion in
\eqref{eq:CD} collapses by dimension counting.  This is the elementary
finite-dimensional core of the non-MF argument.  It does not collapse in an
arbitrary finite von Neumann algebra: proper finite-index subalgebras can be
unitarily conjugate to themselves in one direction inside a larger finite
algebra.  The regular representation of $E$ already supplies an abstract
finite-algebra realization of \eqref{eq:CD}.

\section{Regular-tensor robustness}\label{sec:dilution}

We now prove that the fixed Clifford factor cannot repair the rank blindness
of the adjoint representation.  We use the quotient
$q:E\to\bar G$ from Lemma~\ref{lem:linearquotient}.

\begin{theorem}[Regular-tensor robustness]\label{thm:dilution}
Assume, hypothetically, that $E$ has a trace-preserving hyperlinear
embedding represented by pointwise asymptotically multiplicative maps
\[
 \rho_n:E\longrightarrow U(m_n)
\]
whose normalized traces converge to the regular character.  There are finite quotients
$q_n:\bar G\to Q_n$ such that the maps
\[
 \widetilde\rho_n(x)
 =\rho_n(x)\otimes\lambda_{Q_n}(q_n(q(x)))
\]
define another trace-preserving hyperlinear embedding.  They have the
following properties.
\begin{enumerate}
\item For every $x,y\in E$, the multiplicative defects satisfy the exact
      identity
      \[
       \|\widetilde\rho_n(x)\widetilde\rho_n(y)
          -\widetilde\rho_n(xy)\|_2
       =\|\rho_n(x)\rho_n(y)-\rho_n(xy)\|_2.
      \]
\item The matrices representing $w$ and all eight $d_x$ are stable
      amplifications by identities.  Hence the trace of the negative corner
      and every normalized $*$-moment of the Clifford generators are
      unchanged.
\item On the negative corner the normalized characters of the coordinate
      adjoint representations of both $\Gamma$ and $H$ converge to their
      regular characters.
\item For any fixed generating-set Laplacian of either infinite
      property-(T) group, the normalized ranks of its coordinate
      low-spectral projections tend to zero below the Kazhdan gap.
\item The distinguished unit vectors represented by $d$ remain
      $H$-invariant in the ultraproduct and retain the exact Clifford wall
      against a distinct conjugate.
\end{enumerate}
\end{theorem}

\begin{proof}
By residual finiteness, choose quotient maps $q_n:\bar G\to Q_n$ which
separate every nonidentity element in an increasing exhaustive sequence of
finite subsets of the embedded copy of $\Gamma$.  Write $\lambda_{Q_n}$
for the left regular representation.
The second tensor factor is an honest representation, so the normalized
Hilbert--Schmidt defect identity in part~(1) follows by factoring out the
common unitary $\lambda_{Q_n}(q_n(q(xy)))$.  Since $q(w)=q(c)=1$, and hence
$q(d_x)=1$ for every $x\in\Omega$, the matrices in part~(2) are merely
tensored with an identity.  Functional calculus for the $w$-representative
shows the same for the negative-corner projection.

The new maps still have the regular limiting trace.  Indeed, for $x\ne1$
the first normalized trace already tends to zero, while the second factor
has modulus at most one.  They also remain separating; equivalently, their
limiting trace is the regular character, so the induced ultraproduct
homomorphism is injective.

Let $p_n$ be projection representatives of the negative corner, of ranks
$r_n$, and let $U_{g,n}$ denote compressed unitary representatives.  For
$g\in\Gamma$, the normalized character of the adjoint matrix on the new
negative corner is
\[
 \chi_n(g)
 =\left|\tr_{r_n}(U_{g,n})\right|^2
  \left|\tr_{|Q_n|}(\lambda_{Q_n}(q_n(q(g))))\right|^2.
\]
The second factor equals zero whenever the finite-quotient image of $g$ is
nontrivial.  Thus $\chi_n(e)=1$ and $\chi_n(g)\to0$ for every fixed
$g\ne e$.  The marked adjoint representations therefore converge in
moments to the regular representation of $\Gamma$.  The same is true after
restriction to $H$, because the quotients separate increasing finite subsets
of $\Gamma$ and hence of $H$.

Because $\Gamma$ and $H\cong\Gamma$ are infinite, their regular
representations have no invariant vectors.  Property (T) gives a spectral
gap above zero for either regular Laplacian.  Moment convergence, followed
by polynomial approximation to a spectral cutoff (or the robust
spectral-gap formulation of \cite[Section~7]{DogonVigdorovich}), shows that
the normalized traces of the coordinate spectral projections below any
fixed cutoff strictly inside the respective gap tend to the mass of the
limiting spectral measure at zero, namely zero.  This also follows directly
from Proposition~\ref{prop:cornerregular}; the tensor construction records
that the conclusion is stable while every Clifford statistic is held fixed.

On the other hand, the group relations still say that $d$ centralizes $H$.
Its coordinate representatives have normalized $2$-norm one, and in the
ultraproduct they are $H$-fixed under conjugation.  The relations in
Theorem~\ref{thm:spin} are unchanged, so $d$ remains at squared
normalized $2$-distance $2$ from any distinct Clifford conjugate.  The
distinguished vector has therefore survived entirely inside an adjoint
sector whose normalized coordinate dimension can vanish.  This proves all
five assertions.
\end{proof}

\begin{corollary}[Ordinary adjoint-rank no-go]\label{cor:ranknogo}
Assuming a trace-preserving embedding of $L(E)$, both the $\Gamma$- and
$H$-low-spectral spaces in the negative corner have normalized ranks tending
to zero, while Corollary~\ref{cor:lampzerorank} places the marked lamp
asymptotically inside the latter with an order-one wall.  These conclusions
are unchanged by the regular tensor operation of
Theorem~\ref{thm:dilution}.  Therefore the fixed $M_{16}$ packet and the
ordinary normalized ranks alone cannot imply a positive-density
$H$-fixed excess, nor can they justify finite-dimensional dimension
counting for the inclusion $C\subseteq D$.
\end{corollary}

\begin{proof}
The rank limits and the marked vector are Proposition~\ref{prop:cornerregular}
and Corollary~\ref{cor:lampzerorank}.  Parts~(2)--(5) of
Theorem~\ref{thm:dilution} preserve all of these data under arbitrarily large
finite regular amplifications.  Hence no positive lower bound proportional
to the ambient adjoint dimension can follow from those data.
\end{proof}

\begin{remark}[The tensor does not sacrifice faithfulness]
The dilution in Theorem~\ref{thm:dilution} is not a counterexample built
from a nonfaithful quotient trace.  It starts with a hypothetical
trace-preserving embedding of $L(E)$ and leaves its regular character
unchanged.  It is therefore a stopping theorem for arguments intended to
use faithfulness only through ordinary group traces.
\end{remark}

\begin{remark}[Why the coefficient algebra does not help automatically]
After exact Clifford disintegration, the commutant
$1\otimes M_k(\C)$ occupies a fixed fraction $1/256$ of the adjoint matrix
space.  But multiplication by the lamp does not make all of this coefficient
space $H$-fixed.  Abstractly, its $H$-fixed part is the commutant of the
inverse-projective coefficient action $v|_H$.  Whenever that action is
recovered compatibly at coordinates, its fixed vectors lie in the full
$H$-low-spectral space, whose ordinary normalized rank vanishes by
Proposition~\ref{prop:cornerregular}.  Theorem~\ref{thm:dilution} shows that
finite regular tensoring cannot change this conclusion while preserving the
entire Clifford packet.  Any successful dimension argument must therefore
divide out the irrelevant regular tensor or use a center-valued dimension
over a canonically recovered coefficient algebra.
\end{remark}

\section{The exact remaining gate}

The spin reduction leaves a concrete simultaneous-coordinate problem.  The
following definition isolates the exact sufficient input without assuming
it.

\begin{definition}[Compatible exact coordinate recovery]
\label{def:recovery}
Let $\theta:L(E)\to\mathcal M_\omega$ be a trace-preserving embedding and
work in its negative corner, represented as
$\prod_\omega(M_{r_n}(\C),\tr_{r_n})$.  Put
\[
 C=\theta(\Gamma)'\cap p\mathcal M_\omega p,
 \qquad
 D=\theta(H)'\cap p\mathcal M_\omega p.
\]
We say that $(C,D)$ has \emph{compatible exact coordinate recovery} if
there are unital $*$-subalgebras $C_n,D_n\subseteq M_{r_n}(\C)$ and unitary
lifts $U_{t,n}$ of $p\theta(t)$ such that, as concrete subalgebras of the
corner ultraproduct,
\[
 C=\prod_\omega C_n,
 \qquad D=\prod_\omega D_n,
\]
and, for an $\omega$-large set of indices,
\[
 C_n\subseteq D_n,
 \qquad D_n=U_{t,n}C_nU_{t,n}^*,
\]
\end{definition}

\begin{proposition}[Compatible recovery closes the candidate]
\label{prop:recoverycloses}
No trace-preserving embedding of $L(E)$ can have compatible exact coordinate
recovery in the sense of Definition~\ref{def:recovery}.
\end{proposition}

\begin{proof}
For every relevant $n$, unitary conjugacy gives
$\dim_\C C_n=\dim_\C D_n$.  Since $C_n\subseteq D_n$ and both are
finite-dimensional vector spaces, $C_n=D_n$.  Taking tracial ultraproducts
gives $C=D$.

But $d=tct^{-1}$ centralizes $H$, so $p\theta(d)\in D$.  Equality $C=D$
would put $p\theta(d)$ in the $\Gamma$-commutant.  Since $a\in\Gamma$, this
would give
\[
 p\theta(ada^{-1})=p\theta(d)
 \quad\text{and hence}\quad
 p\theta(w)=p.
\]
The definition of the negative corner instead gives $p\theta(w)=-p$.
As $p\ne0$, this is a contradiction.
\end{proof}

Alekseev--Thom Open Problem~6.2 asks whether the commutant of a Kazhdan
representation in a tracial matrix ultraproduct can be recovered as an
ultraproduct of finite-dimensional algebras, and whether these can be taken
as coordinate centralizers \cite{AlekseevThom}.  The present HNN argument
needs, in addition, simultaneous recovery of both commutants, compatibility
with the same coordinate lifts of $t$, and exact nesting (or a perturbative
replacement strong enough to imply it).  Neither statement is presently
available.  Proposition~\ref{prop:recoverycloses} is therefore a conditional
closure theorem, not an invocation of the open input.

Theorem~\ref{thm:dilution} shows that replacing this gate by ordinary
low-spectral-space rank is impossible.  A viable replacement must be one of
the following genuinely relative statements:
\begin{enumerate}
\item simultaneous coordinate recovery of $C$ and $D$ with HNN-compatible
      nesting;
\item a center-valued or module dimension over the recovered Clifford
      coefficient algebra which is invariant under the regular tensor; or
\item a canonical finite-multiplicity leakage inequality which sees a
      bounded unitary of $D\setminus C$ even when its ambient adjoint sector
      has trace zero.
\end{enumerate}
Each would be a new matrix-specific theorem.  None follows from abstract
finite-von-Neumann-algebra finiteness, from the fixed $M_{16}$ factor, or
from property (T) spectral gap alone.

\section{Conclusion}

The eight-lamp calculation is exact: a hypothetical hyperlinear embedding
of $E$ contains, on a corner of trace $1/2$, a forced
$M_{16}(\C)$ spin factor and an inverse-projective coefficient system.
The spin factor admits exact coordinate lifts
$M_{16}\otimes1_{k_n}\subseteq M_{16k_n}$, and the coefficient commutants
form a proper finite-index one-sided HNN inclusion.

The new boundary theorem is equally exact.  The negative-corner adjoint
characters of $\Gamma$ and $H$ are already regular, so their normalized
low-spectral ranks vanish.  A finite regular tensor leaves the entire
Clifford package and faithful group trace unchanged and shows that this rank
blindness is robust.  Thus the spin factor removes
lamp-gauge ambiguity but does not remove second-order rank blindness.  The
remaining endpoint is a relative finite-multiplicity theorem, not an
ordinary Kazhdan spectral-projection estimate.

Definition~\ref{def:recovery} and
Proposition~\ref{prop:recoverycloses} formalize the endpoint: simultaneous
exact coordinate recovery of the two commutants, compatible with the HNN
unitary, would close the candidate immediately.  This conditional theorem
is deliberately separated from the unavailable recovery input.

Accordingly, this note does not assert the requested nonhyperlinear group.
It proves the strongest currently verified reduction from this marked
candidate and rules out a natural but invalid closure of that reduction.

\clearpage
\begin{thebibliography}{9}

\bibitem{AlekseevThom}
V. Alekseev and A. Thom,
\emph{Centralizers of sofic approximations of Kazhdan groups},
arXiv:2608.05362, 2026.

\bibitem{BHV}
B. Bekka, P. de la Harpe, and A. Valette,
\emph{Kazhdan's property (T)},
New Mathematical Monographs, vol.~11,
Cambridge University Press, Cambridge, 2008.

\bibitem{DogonVigdorovich}
A. Dogon and I. Vigdorovich,
\emph{Hyperlinearity, stability and asymptotic spectral gap of higher rank
lattices},
arXiv:2506.20843v2, 2026.

\bibitem{Malcev}
A. I. Mal'cev,
\emph{On isomorphic matrix representations of infinite groups},
Mat. Sb. (N.S.) \textbf{8(50)} (1940), no.~3, 405--422.

\bibitem{TallerVidick}
A. Taller and T. Vidick,
\emph{Approximating the quantum value of an LCS game is RE-hard},
arXiv:2507.22444v2, 2026.

\end{thebibliography}

\end{document}
